\section{Decision-problem hardness: complete reductions}
\label{sec:annex-p04-decision-hardness}

\begin{annexitem}{annex:p04.dh.stcon}{thm:p04-dh-stcon}{Configuration graphs and logarithmic-space hardness}
\begin{proof}
Reject a malformed graph encoding. For a valid \(N\)-vertex graph, a path
exists iff one exists with at most \(N-1\) edges: remove the segment between
two repeated vertices from any longer witness until none repeat. A
nondeterministic machine stores its current vertex and a counter bounded
by \(N-1\), guesses successive vertices and checks the adjacency bit by
rescanning the explicit matrix. It may accept whenever it reaches \(t\),
including before a move when \(s=t\). Every branch stops at its length
bound, and all stored numbers need \(O(\log(N+1))\) bits. This proves
membership under the stated all-branches-halting convention.

For hardness fix any language \(A\in\mathsf{NL}\) and its machine \(M\).
On an input \(x\) of length \(n\), its configurations consist of the finite
control state, the input-head position, the work-head positions and the
contents of at most \(C\lceil\log(n+2)\rceil\) work cells, for a fixed \(C\).
Encode all these fields in \(b(n)=O(\log(n+2))\) bits, increasing the fixed
constant to allow separators and state/head fields. There are at most
\(2^{b(n)}\), hence polynomially many, candidate configurations. Include
all candidate bit strings, not only reachable ones; malformed configurations
are isolated vertices. Add one additional accepting sink \(a\).
The start vertex is the valid initial configuration; edges join valid
configurations related by one transition of \(M\), and every accepting
configuration has an edge to \(a\).

To test a candidate edge, keep its two configurations in logarithmic
workspace, check their fields and work-space boundary, and compare the
local transition with the input symbol at the stored input-head position.
The machine \(M\) is fixed, so its finite transition relation is constant
data. Enumerate vertex indices and pairs with binary counters and emit the
unary count, designated indices and adjacency matrix one bit at a time.
The output has polynomial length, although it need not fit in workspace.
The translator is deterministic logarithmic-space and outputs a valid
graph on every \(x\), including \(n=0\). There is a path from the initial
configuration to \(a\) iff a branch of \(M\) accepts \(x\). This is a
many-one reduction from the arbitrary \(A\), proving hardness.

We will also use the standard requested-bit implementation of these
translations. To simulate a logarithmic-space machine on a
polynomial-length generated word without storing that word, keep its
input-head index and rerun the output transducer to obtain the requested
bit. Each index and both work areas are \(O(\log(n+2))\). This proves that
logspace reductions compose and that NL membership pulls back through
them; it is not an assumption that the generated graph is resident in memory.
\end{proof}
\end{annexitem}

\begin{annexitem}{annex:p04.dh.complement}{thm:p04-dh-complement}{Inductive counting without a negative reachability oracle}
\begin{proof}
The construction expands the two counting lemmas of
\cite[theorem 1 and lemmas 1--2]{Immerman1988Complement} for the explicit
finite graph convention used here. Reject malformed input first. For
a valid graph, the sets \(R_d\) increase and satisfy
\[
 v\in R_{d+1}\quad\Longleftrightarrow\quad
 \exists u\in R_d\;(u=v\text{ or }(u,v)\in E).
\]
A path certificate of length at most \(d\) is verified by storing its
current vertex, destination and length counter. The length-zero option
is present, so the base set is exactly \(R_0=\{s\}\).

Suppose \(c_d=|R_d|\) is correct. In one scan over \(u=1,\ldots,N\),
an included vertex must supply a verified length-at-most-\(d\) path.
Thus no surviving branch includes a vertex outside \(R_d\), and the
ordered scan prevents counting any vertex twice. If its final inclusion
count equals \(c_d\), it must have included every vertex in \(R_d\).
This is why omission is not an unjustified negative answer: omitting a
reachable vertex makes the count too small. Conversely there exists a
successful branch, choosing inclusion and a witnessing path for each
\(u\in R_d\), and omission otherwise.

For a fixed candidate \(v\), set its flag precisely when an included
\(u\) equals \(v\) or has an edge to \(v\). On every branch passing the
count check that flag is therefore exactly the membership indicator
of \(R_{d+1}\). Repeat this full checked scan for each \(v\); increment
a fresh counter for each true flag. All surviving branches then compute
the same correct \(c_{d+1}\), and successful choices exist for all scans.
Starting at \(c_0=1\) and iterating \(N-1\) times computes \(c_{N-1}\).
There is no recursive subroutine storing older counts or path histories.

Every reachable vertex has a simple witnessing path, so
\(R_{N-1}\) is the full reachable set. Make one last ordered scan with
path certificates of length at most \(N-1\), require exactly \(c_{N-1}\)
inclusions, and accept iff \(t\) was not included. If \(t\) is reachable,
every count-checked branch includes it and rejects. If not, successful
choices include all reachable vertices and accept. For \(N=1\), there
is no update stage and the final scan includes \(s=t\), so the no-path
test correctly rejects.

A fixed number of counters for the stage, candidates, inclusion count,
new count and path length, plus a fixed number of vertex indices and
flags, fit in \(O(\log(N+1))\) work cells. Path guessing and every loop
are bounded by \(N\); every branch halts, including branches that fail
a check. Rescanning the input to access matrix bits is logarithmic-space.
Thus \(\mathsf{NONREACH}\in\mathsf{NL}\).

For any \(A\in\mathsf{NL}\), the previous proof supplies a valid
configuration graph \(g(x)\) for every binary word \(x\), with
\(x\in A\) iff \(g(x)\in\mathsf{STCON}\). Therefore
\(x\notin A\) iff \(g(x)\in\mathsf{NONREACH}\).
The requested-bit simulation shows \(\overline A\in\mathsf{NL}\).
Applying the same statement twice establishes
\(\mathsf{NL}=\mathsf{coNL}\) as classes over the whole binary input space.

Finally, if \(A\in\mathsf{NL}\), then \(\overline A\in\mathsf{NL}\);
apply the configuration reduction to \(\overline A\).
Its always-valid output satisfies \(x\in A\) iff it is a NONREACH yes
instance. This proves NL-hardness of NONREACH. The argument never
mistakes valid-encoding nonreachability for the whole complement of STCON.
\end{proof}
\end{annexitem}

\begin{annexitem}{annex:p04.dh.two-sat}{thm:p04-dh-two-sat}{The SCC criterion, the NL upper bound, and the lower reduction}
\begin{proof}
Reject malformed input and handle an empty clause separately. On a valid
input with no empty clause, form the implication relation on the \(2r\)
signed literals: \(L\lor K\) contributes
\(\overline L\to K\) and \(\overline K\to L\); a unit repeats its literal.
These edges can be queried in logarithmic space by scanning the clauses.

Here is a full justification of the criterion used by
\cref{thm:p02-two-sat}. Every satisfying assignment respects each edge,
and hence each directed path. If \(p\) and \(\neg p\) lie in the same
SCC, either truth choice for \(p\) forces its opposite, so satisfaction
is impossible. Conversely suppose no SCC contains an opposite pair.
Edge reversal followed by literal complementation preserves the graph:
an edge \(L\to K\) always has the companion
\(\overline K\to\overline L\). Thus complementation maps each component
\(C\) to a distinct component \(\overline C\), reversing the condensation
DAG's edges. Choose a topological rank \(h(C)\), increasing on edges
between distinct components; the finite DAG has such a rank by repeatedly
removing a vertex of indegree zero. Define a literal to be true iff
\[
 h(C_L)>h(C_{\overline L}).
\]
Exactly one literal in each opposite pair is true. An edge \(L\to K\)
with \(L\) true and \(K\) false would imply
\[
 h(C_{\overline L})<h(C_L)\leq h(C_K)
 <h(C_{\overline K})\leq h(C_{\overline L}),
\]
a contradiction, using the companion edge for the last inequality.
Hence every edge is respected. If a clause \(L\lor K\) were false,
\(\overline L\) would be true and its edge to \(K\) would force \(K\)
true, another contradiction. This constructs a satisfying assignment.
For \(r=0\), a valid input without an empty clause is the empty clause
list and is satisfied by the empty assignment.

On well-formed inputs, unsatisfiability is therefore certified by an
empty clause or a guessed \(i\) and two bounded paths
\(p_i\leadsto\neg p_i\leadsto p_i\). Each path needs at most \(2r-1\)
edges. The nondeterministic machine stores only a fixed number of
\(O(\log(r+1))\)-bit indices and counters; adjacency is queried directly
from the input. Syntax checking is also logarithmic-space. Consequently
the whole complement of 2SAT, which includes malformed inputs as well
as well-formed unsatisfiable inputs, belongs to NL. Complementation
proved above gives \(\mathsf{2SAT}\in\mathsf{NL}\).

For hardness translate a valid graph to the declared \(r=N\) variables
and emit
\[
 p_s,\quad \neg p_t,\quad
 (\neg p_u\lor p_v)\quad\text{for every }(u,v)\in E.
\]
If \(s\leadsto t\), the positive initial unit and the edge implications
force \(p_t\), contrary to its negative unit. If \(t\) is not reachable,
assign exactly the reachable vertices true. A true antecedent edge has
a reachable successor, so every implication is satisfied; the two
units also hold. Therefore the emitted formula is satisfiable exactly
for NONREACH yes instances. Row/column counters and binary variable
indices need logarithmic workspace; output length is polynomial.
Malformed graph inputs map to the fixed valid contradiction
\(p_1\land\neg p_1\). Thus this is a total logspace many-one reduction.
The case \(s=t\) produces contradictory units, as required even for
an edgeless graph. NL-hardness of NONREACH now proves the claimed lower
bound.

The same graph formula is unsatisfiable exactly for STCON yes instances.
Map a malformed STCON input instead to the fixed satisfiable one-variable
empty clause list. Its complement language on valid formula encodings
therefore has the matching NL hardness and the already established
upper bound.

For arbitrary finite atom strings, enumerate occurrences by input
position. Compare two names character by character using their position
and length counters. To find whether an occurrence is first, compare
it with earlier occurrences; to compute its normalized index, count
earlier first occurrences. The fixed-depth nested scans use logarithmic
workspace and polynomially many comparisons. Count distinct names to
emit the unary declaration and translate each literal by rescanning.
Malformed strings map to a fixed contradiction. The reverse translation
prints the short indexed atom names and discards unused declarations.
These details establish the stated encoding equivalence, rather than
assuming that very long names can reside in logarithmic memory.
\end{proof}
\end{annexitem}

\begin{annexitem}{annex:p04.dh.computation}{thm:p04-dh-computation-validity}{Actual transition simulation and a free-term countermodel}
\begin{proof}
First verify the stack encoding against the tape operation. For a nonempty
right list \(\mathsf a_a(v)\), its head is the scanned symbol \(a\);
for an empty right list the scanned symbol is blank and its tail is
empty. A right move writes \(b\), pushes \(b\) onto the left list and
exposes the old right tail. A stationary move replaces the scanned
symbol by \(b\). A left move with left list \(\mathsf a_d(u)\) exposes
\(d\) as the new scanned symbol and leaves the written \(b\) immediately
to its right, giving \(\mathsf a_d(\mathsf a_b(v))\).
If the left list was empty, the exposed symbol is blank and the new
right list is \(\mathsf a_{\sqcup}(\mathsf a_b(v))\).
When the right list was empty, replacing its pattern and tail by
\(\mathsf{nil}\) gives exactly the same operation on the unrepresented
blank region. Explicit blank entries are retained; no normalization
by deleting them is presumed.

Starting from \((q_0,\mathsf{nil},[w])\), these updates inductively
represent every step of the actual machine run relative to its current
head. Absolute head position is not needed by a transition table depending
only on state and scanned symbol. Different list representations of the
same all-blank region cause no problem: the update from each represented
configuration has exactly its actual successor representation. The
injective alphabet and state assignments make the relevant head pattern
unique. Each list remains a finite well-formed list even after any finite
number of moves into an initially blank region.

Suppose the run reaches \(q_{\mathrm a}\) after \(k\) steps. In any
structure satisfying \(T_M\land I_{M,w}\), the initial ground Reach atom
is true. At each step, instantiate the appropriate universally closed
transition clause with the actual two ground list tails. Modus ponens
gives its successor Reach atom. Induction over these \(k\) metatheoretic
steps gives a Reach atom with accepting constructor, whose lists witness
\(H_M\). This proves \(\Phi_{M,w}\) valid in every structure, without
assuming that its functions are injective or that its Reach predicate
contains only actual configurations. It also constructs a finite
natural-deduction derivation of the implication from the closed
transition axioms and initial atom. Immediate acceptance uses \(k=0\).

Conversely assume there is no accepting state in the run, whether it is
infinite or ends at a missing transition. Let the domain \(D\) be all
finite ground terms over \(\Sigma_{\mathrm H}\). It is nonempty because
it contains \(\mathsf{nil}\). Interpret every function as the corresponding
term constructor, including unused symbols, and the constant as its
own ground term. Constructors are injective, have disjoint images, and
none of their values is the constant. If identity occurs in the ambient
language, interpret it as literal term identity.

Interpret Reach as the set of exactly the configuration terms produced
by the represented run. The initial atom is therefore true. Consider
any assignment to the variables of any transition clause. If its
antecedent is false, the clause holds. If true, its ground term is an
actual run configuration. Freeness forces its state constructor and
list pattern to match exactly the represented state and lists. A matching
transition clause is precisely an applicable table entry; the verified
update gives the next represented run term, which also belongs to Reach.
Thus every universal clause is true, including assignments of arbitrary
non-list ground terms: those cannot make a falsely typed antecedent
equal an actual run term. This is why no typing axioms are missing.

No term headed by the distinct accepting-state constructor belongs to
Reach, since the run does not accept. Thus every choice of \(u,v\in D\)
makes the matrix of \(H_M\) false. The structure satisfies
\(T_M\land I_{M,w}\) and falsifies \(H_M\), so it is a concrete countermodel
to \(\Phi_{M,w}\). It may be infinite, as the unrestricted target permits.

Finally the compiler is total on all strings: parse a finite machine/word
description, perform the stated syntactic checks, and otherwise return
the fixed closed contradiction
\(\mathsf{Reach}(\mathsf{nil})\land\neg\mathsf{Reach}(\mathsf{nil})\).
For a valid pair, enumerate its finitely many transitions and symbol
cases and form the finite sentence. A left transition emits at most
\(|\Gamma|+1\) clauses for each of at most two right patterns; other
moves emit at most two. Hence the stated clause bound and polynomial
written size hold in the explicit-list encoding. No step runs the
machine to determine whether compilation should finish.
\end{proof}
\end{annexitem}

\begin{annexitem}{annex:p04.dh.validity}{thm:p04-dh-validity-complete}{The reduction from every c.e. set}
\begin{proof}
For the upper bound, enumerate finite derivations in the effective
first-order calculus, check that they have no undischarged assumptions,
and output their closed conclusions in \(\Sigma_{\mathrm H}\).
Decidable arities and syntax let us discard other strings. Soundness
ensures every output is valid; the existing first-order completeness
theorem ensures every valid closed sentence is eventually output.
This is exactly the enumerability argument at
\cref{thm:fol-validity-enumerable}, specialized to the declared countable
effective signature. On a requested code, syntactic invalidity is
rejected and valid syntax is semidecided by the proof enumeration.
The empty domain is never used.

For the lower bound take any c.e. \(A\subseteq\N\). Its actual defining
property in \cref{sec:turing-machines-03} is that \(A=\operatorname{dom}(f)\)
for some partial computable \(f\). Unpack
\cref{def:partial-computable}: there is a finite machine \(D_A\) which
on each canonical \(\operatorname{bin}(x)\) accepts with the proper
numerical output when \(f(x)\) is defined and has an infinite computation
when it is undefined. Regard that machine as a recognizer \(M_A\),
ignoring its numerical output after acceptance. Then it enters its
accepting state exactly for \(x\in A\). This is a direct proof of the
recognizer conversion needed for this reduction; it does not depend
on the separately sketched equivalence of all c.e. presentations or on
bounded arithmetic representability. Multiple accepting halts can be
redirected to one fresh halting accepting state, and finitely many
alphabet names can be indexed injectively, by effective finite-table
transformations. The two-way tape and binary input convention agree
with the stack compilation.

With this fixed finite \(M_A\), define the total computable map
\[
 f_A(x)=\operatorname{code}
        \bigl(\Phi_{M_A,\operatorname{bin}(x)}\bigr).
\]
Computing the finite word \(\operatorname{bin}(x)\), emitting its nested
list term and the already fixed transition clauses, and applying the
effective sentence coding all terminate. This computation does not
wait for a run of \(M_A\). By the preceding lemma,
\[
 x\in A
 \quad\Longleftrightarrow\quad M_A\text{ accepts }\operatorname{bin}(x)
 \quad\Longleftrightarrow\quad
 f_A(x)\in\mathsf{VALID}_{\mathrm H}.
\]
Thus every c.e. set reduces by \(\leq_m\) to the target, establishing
c.e.-hardness, and the upper bound gives completeness.

The reduction is uniform as a compiler in a supplied machine description
and word. The choice of a fixed recognizer for each fixed c.e. set is the
usual quantifier order for many-one hardness, not a claim that arbitrary
descriptions of sets can be recognized without their machine data.
For \(A=\varnothing\), take a machine that never accepts; every resulting
sentence has the free-term countermodel. For \(A=\N\), take a machine
accepting every canonical binary input; every resulting sentence is valid.
The number zero is represented by the canonical word \(0\), not an
empty numeral. The separate empty word is nevertheless handled by the
compilation lemma. A valid empty transition table accepts only if its
initial state is already accepting; otherwise its initial Reach atom
and absent accepting constructor give the same countermodel.

General first-order validity contains these effectively coded instances:
rename each indexed signature symbol into the ordinary finite formula
signature and apply the fixed computable recoding. Hence the reduction
also proves its c.e.-hardness. General proof enumeration supplies its
c.e. upper bound. No further assertion for every individual fixed
finite signature follows from this argument.
\end{proof}
\end{annexitem}

\begin{annexitem}{annex:p04.dh.relational}{thm:p04-dh-relational-validity}{Effective function elimination with both model directions}
\begin{proof}
It suffices to transform the finite-symbol sentence \(\Phi_{M,w}\).
Choose a variable \(z_0\) absent from it. For every used \(k\)-ary
function \(f\), where \(k=1\) or \(2\), introduce a fresh \((k+1)\)-ary
predicate \(G_f\). Form the two closed axioms
\[
 \forall x_1\cdots\forall x_k\,\exists y\,G_f(\vec x,y),
 \qquad
 \forall\vec x\,\forall y\,\forall y'\,
 ((G_f(\vec x,y)\land G_f(\vec x,y'))\to y=y').
\]
Their finite conjunction is \(J\); it contains no \(z_0\).
It asserts genuine total functional graphs, not merely that some
relations have been named.

For each term \(t\) define a formula \(V_t(y)\), with every auxiliary
bound variable freshly chosen outside all surrounding variables:
\[
\begin{aligned}
 V_x(y)&=(y=x),\\
 V_{\mathsf{nil}}(y)&=(y=z_0),\\
 V_{f(t_1,\ldots,t_k)}(y)&=
  \exists y_1\cdots\exists y_k\,
  \left(\bigwedge_{i=1}^k V_{t_i}(y_i)
                 \land G_f(y_1,\ldots,y_k,y)\right).
\end{aligned}
\]
Replace an atom \(\mathsf{Reach}(t)\) by
\(\exists y\,(V_t(y)\land\mathsf{Reach}(y))\).
Retain Boolean connectives and quantifiers recursively, choosing generated
variable names to avoid capture. Fresh names can be selected effectively
from the infinite variable list by skipping its finite forbidden set.
Write \(\widehat\Phi(z_0)\) for the resulting formula, and output
\[
 \Psi=\forall z_0\,(J\to\widehat\Phi(z_0)).
\]
All graph axioms were closed before this implication was formed.

In any relational structure satisfying \(J\), totality and uniqueness
give one actual function for each \(G_f\). For an assignment to \(z_0\),
interpret the old constant as that element. Induction on terms proves
that \(V_t(y)\) holds for exactly the value of \(t\): the variable and
constant cases are identities; the functional case first obtains the
unique child values and then the unique graph output. Induction on
formulas now gives equivalence of \(\widehat\Phi\) and \(\Phi\).
The atomic case follows from that unique value; negation and the Boolean
connectives preserve equivalence; a quantifier varies the same original
variable over the same nonempty domain, with all generated variables
fresh. This establishes the full truth correspondence, not just a
one-way satisfiability implication.

If \(\Phi\) is valid, any relational structure and any \(z_0\) assignment
either fail \(J\), making the implication true, or induce the function
structure above, in which \(\Phi\) and hence \(\widehat\Phi\) are true.
Thus \(\Psi\) is valid. Conversely a countermodel to \(\Phi\) supplies
a relational model by taking each function's actual graph and assigning
\(z_0\) its constant value. It satisfies \(J\) and falsifies
\(\widehat\Phi\), so falsifies the universal sentence \(\Psi\).
Both model directions use arbitrary nonempty structures.

The finite term recursion, fresh-name choices and graph-axiom emission
are total computable operations. They produce only the unary Reach
predicate, binary graphs for unary functions and ternary graphs for
binary functions, together with logical identity; there are no remaining
function or constant symbols. Composing this translation with each
\(f_A\) therefore proves c.e.-hardness of the general relational validity
language with its finite signature encoded in the input. Enumerating
effective first-order proofs and retaining relational closed conclusions
gives the c.e. upper bound. If a malformed compiler input needs an image,
the fixed relational contradiction
\(\forall x\,(\mathsf{Reach}(x)\land\neg\mathsf{Reach}(x))\) is a closed
nonvalid sentence on every nonempty domain. The reduction is total,
and no uniform claim for all fixed finite relational signatures is made.
\end{proof}
\end{annexitem}
